\section{An effective Catalan asymptotic}
\label{sec:analytic-catalan}
The links to \emph{Analytic combinatorics} and \emph{Catalan number}
lead from an exact counting recurrence to an estimate of its coefficients.
In Section~\ref{sec:enumeration}, $C=1+xC^2$ gives
$(1-2xC)^2=1-4x$. Thus $1-2xC$ is the formal square root with
constant coefficient $1$: coefficient $j>0$ of a square of this form
equals twice coefficient $j$ of the root plus a sum over smaller indices.
This proves uniqueness by triangular comparison. We now obtain an
asymptotic with a numerical error at every positive index.

\begin{theorem}[Wallis bounds and relative error]\label{result14}
Set $a_n=4^n/\binom{2n}{n}$ for $n\ge1$. With the constant $\pi$
specified by rational approximants below,
\[
 L_n=\frac{2a_n^2}{2n+1}\ \le\ \pi\ \le\
 U_n=\frac{a_n^2}{n},\qquad U_n-L_n\le\frac4{2n+1}.
\]
Consequently the two normalized counts satisfy
\[
\begin{split}
 0&\le 1-\frac{\sqrt{\pi n}}{4^n}\binom{2n}{n}\le\frac1{4n},\\
 0&\le 1-\frac{\sqrt{\pi}\,n^{3/2}C_n}{4^n}
       \le\frac5{4(n+1)}.
\end{split}
\]
Thus $C_n\sim 4^n/(\sqrt{\pi}\,n^{3/2})$ with, for example, the
primitive recursive relative-error modulus $n\ge 2^{k+1}$ for tolerance
$2^{-k}$.
\end{theorem}

\paragraph{Rational meaning of the integral notation.}
For a rational function $f=P/Q$ used below, the displayed denominator
satisfies $Q(t)\ge1$ on $0\le t\le1$. Write
$S_M(f)=M^{-1}\sum_{i=0}^{M-1} f(i/M)$ for $M\ge1$.
The difference $P(t+h)Q(t)-P(t)Q(t+h)$ has zero constant coefficient
in $h$. Divide by $h$ in the polynomial ring, and let $K$ be the
sum of the absolute coefficients of that quotient, rounded up.
For $0\le t\le t+h\le1$ the denominators of $f(t+h)-f(t)$ have
product at least one, so $|f(t+h)-f(t)|\le Kh$.
Comparing $S_M$ with the refinement $S_{MN}$ on each original cell
gives an error at most $K/M$. Comparing with $S_N$ in the same way gives
$|S_M-S_N|\le K/M+K/N$. The notation $\int_0^1 f$ stands for this
explicit rational approximation procedure. At precision $2^{-k}$ one
may take $M=(K+1)2^{k+2}$. Sums, rational scalar multiplication, and
order inequalities on a common mesh follow directly from the finite sum.

One other calculation suffices in place of a general integration theorem.
For $F=P/Q$ form the polynomial
\[
 E(t,h)=(P(t+h)Q(t)-P(t)Q(t+h))Q(t)
       -h(P'Q-PQ')(t)Q(t+h).
\]
Its constant and linear terms in $h$ vanish by the binomial expansion.
Let $J(F)$ be the ceiling of the sum of absolute coefficients of
$E/h^2$. Dividing by $Q(t+h)Q(t)^2\ge1$ gives
$|F(t+h)-F(t)-hF'(t)|\le J(F)h^2$, with
$F'=(P'Q-PQ')/Q^2$ computed formally. Telescoping at $h=1/M$ proves
\[
 |S_M(F')-(F(1)-F(0))|\le J(F)/M.                  \tag{*}
\]
These constants are primitive recursive functions of the coefficient arrays.

\paragraph{Identifying the constant.}
Define $\pi$ by $\pi_M=4S_M(1/(1+t^2))$. To identify the name independently,
the finite geometric identity and $(*)$ for
$F=t^{j+1}/(j+1)$ give Gregory's classical characterization
\[
 \left|\pi-4\sum_{j=0}^{b-1}\frac{(-1)^j}{2j+1}\right|
       \le\frac4{2b+1}.
\]
Indeed the pointwise remainder before integration has absolute value
$t^{2b}/(1+t^2)\le t^{2b}$. There is also a geometric identification.
The parametrization $v(t)=((1-t^2)/(1+t^2),2t/(1+t^2))$ traverses the
unit quarter circle for $0\le t\le1$; its inverse is $t=y/(1+x)$.
Direct squaring shows the chord from $t$ to $u\ge t$ has length
$2(u-t)/\sqrt{(1+t^2)(1+u^2)}$. This lies between the right- and
left-endpoint rectangle contributions for the decreasing function
$2/(1+t^2)$. On a uniform mesh those two totals differ by
$1/M$ by telescoping. Thus the polygonal quarter-arc length converges
to $\pi/2$. For any rational partition refine it by the uniform mesh;
the coarse chord sum can only increase under refinement by the triangle
inequality, and the refined upper and lower endpoint totals differ by
at most $1/M$ (sum the drops multiplied by the largest mesh length).
This also identifies the constant with the usual circle circumference
divided by its diameter. The arithmetic proofs can use the rational
$\pi_M$ throughout.

\begin{proof}[Proof of the bounds]
Put $c=(1-t^2)/(1+t^2)$, $s=2t/(1+t^2)$, $w=2/(1+t^2)$, and
$I_j=\int_0^1 wc^j$. The identities $c^2+s^2=1$, $s'=wc$ and
$c'=-ws$ follow by rational polynomial multiplication. We have
$I_0=\pi/2$, $I_1=1$ from $s(1)-s(0)=1$ and $(*)$. For $j\ge2$
differentiate the product $sc^{j-1}$ formally:
\[
 (sc^{j-1})'=jwc^j-(j-1)wc^{j-2}.
\]
The endpoint values vanish, giving $j I_j=(j-1)I_{j-2}$. Since
$0\le c\le1$, the $I_j$ decrease. The recurrence gives
\[
 I_{2n}=\frac{\pi}{2a_n},\qquad
 I_{2n+1}=\frac{a_n}{2n+1},\qquad I_{2n-1}=\frac{a_n}{2n}.
\]
For example the even product multiplier is
$\prod_{j=1}^n(2j-1)/(2j)=\binom{2n}{n}/4^n$, by separating the
factors of $(2n)!$. Comparison of the three integrals gives the enclosing
interval. Also $a_{n+1}/a_n=2(n+1)/(2n+1)$, so
$L_{n+1}/L_n=(2n+2)^2/((2n+1)(2n+3))>1$ and
$U_{n+1}/U_n=4n(n+1)/(2n+1)^2<1$. The width is
$U_n/(2n+1)\le U_1/(2n+1)=4/(2n+1)$.

For an explicit iPRA induction in this argument fix the target index
$m=2n+1$ and one mesh $M$. Replace $I_j$ by $S_M(wc^j)$.
Compare these rational numbers with a rational table $u_j$ obeying the
exact recurrence with $u_0=\pi_M/2$ and $u_1=1$. Let $J$ be the largest
of $1,J(s)$ and the $J(sc^{j-1})$ for $2\le j\le m$.
Equation $(*)$ bounds the error in each recurrence, after division by
$j$, by $J/M$. Since $(j-1)/j\le1$, induction on the finite row
gives $|u_j-S_M(wc^j)|\le(m+1)J/M$ for all $j\le m$.
Thus $\pi_M$ satisfies both Wallis bounds to within
$4a_n(m+1)J/M$. For any requested precision the product in this numerator
specifies how large to take $M$. This proof only inducts on a bounded
Boolean comparison of rational tables.

Finally let $R_n=\sqrt{\pi n}/a_n$. Squaring and using the bounds yields
$1/(1+1/(2n))\le R_n^2\le1$.
For rational $x\ge0$ the identity
$1-1/\sqrt{1+x}=x/(\sqrt{1+x}(1+\sqrt{1+x}))\le x/2$
gives $1-R_n\le1/(4n)$. The Catalan ratio equals
$nR_n/(n+1)$ by the exact counting formula; its deficit is at most
$1/(n+1)+(n/(n+1))/(4n)=5/(4(n+1))$.
One may eliminate the square roots in the internal assertions by squaring
their nonnegative endpoints. Alternatively a square root of a bounded
nonnegative rational is approximated by bisecting a rational interval
with a square comparison, halving its width at each step.
\end{proof}

\paragraph{Lean targets.}
Define rational coefficient arrays, polynomial translation
$P(t+h)$ by binomial expansion, and the nonnegative integers
\texttt{lipBound}, \texttt{taylorBound} from the coefficient sums above.
Prove \texttt{sumRefine} on the rectangular index set $[M]\times[N]$
and \texttt{telescopeDerivative} by summing the polynomial remainder
inequality. Implement \texttt{wallisRow} and \texttt{riemannRow} for a
supplied $m,M$. The central lemma bounds the largest difference on every
completed prefix by its length times $J/M$. Comparison of entries
$2n-1,2n,2n+1$ and positive cross multiplication then yields the
finite precision theorem. These are outline targets.
